\chapter{面向矿井巷道定位的高精地图瓦片化组织与动态加载方法}

\section{引言}

第三章围绕多源结构先验约束下的高精地图构建与优化方法展开研究，获得了可用于矿井巷道定位的高精先验地图；第四章在此基础上，研究了融合多源结构先验信息的定位方法，实现了复杂矿井场景下的连续定位。然而，在实际应用过程中，随着矿井巷道规模不断扩大，若仍采用整图一次性加载方式，不仅会带来较大的内存占用和计算负担，还会增加匹配过程时延，影响定位系统的实时性与工程可部署性。特别是在巷道分支较多、空间范围较大的场景下，整图加载方式难以满足在线定位对高效地图调用的需求。

针对上述问题，本章围绕“高精地图如何高效服务于定位”这一目标，提出面向矿井巷道定位的高精地图瓦片化组织与动态加载方法。该方法在不改变正式高精地图构建方式的前提下，将离线优化后的全局高精地图划分为多个局部瓦片，并构建相应的元数据索引；在定位运行阶段，根据当前位姿估计结果动态选择所需瓦片，构建局部匹配地图，从而降低整图加载造成的存储与计算开销，提高大尺度矿井环境下定位系统的实时性与稳定性。

需要说明的是，本章的研究重点不是重新构建高精地图，而是解决高精地图在定位系统中的组织、调度与调用问题，其作用在于为第四章提出的定位方法提供高效、稳定、可扩展的先验地图支撑。整体流程如图~\ref{fig:ch5_framework} 所示。

\begin{figure}[htbp]
    \centering
    \fbox{
        \parbox{0.92\linewidth}{
            \centering
            图5.1 面向矿井巷道定位的高精地图瓦片化组织与动态加载总体框架图占位图\\[4pt]
            可按“离线地图构建与切分、元数据生成、定位节点、地图调度模块、地图服务模块、局部地图匹配定位模块”的流程进行绘制。
        }
    }
    \caption{面向矿井巷道定位的高精地图瓦片化组织与动态加载总体框架图}
    \label{fig:ch5_framework}
\end{figure}

\section{定位系统总体架构}

为实现高精地图对矿井巷道定位的有效支撑，本文构建了由地图预处理模块、地图索引模块、动态调度模块、地图服务模块以及定位匹配模块组成的定位系统总体架构。如图~\ref{fig:ch5_framework} 所示，系统主要包含离线与在线两部分。

离线部分主要完成高精地图瓦片化组织与元数据生成。首先，将第三章得到的全局高精地图输入地图分区工具，根据预设的空间分辨率对整幅地图进行规则瓦片划分，生成多个局部点云瓦片；随后，为每个瓦片建立对应的索引信息，包括瓦片文件名、空间起始坐标以及分区分辨率等，形成地图元数据文件。该过程不改变地图本身的几何精度，仅对高精地图进行结构化组织，为后续在线调用提供基础。

在线部分主要完成动态地图调用与定位计算。定位节点实时接收激光点云数据，并结合前一时刻位姿估计结果，判断当前位置所属地图区域；地图调度模块根据当前位置和邻域范围选择所需瓦片集合，并通过地图请求消息发送给地图服务模块；地图服务模块读取对应瓦片并拼接生成局部地图，再将局部地图发布给匹配定位模块；最终，定位模块利用当前激光帧与局部地图进行配准，输出实时位姿估计结果。

该架构将“地图组织”与“地图调用”解耦，使定位系统能够根据空间位置按需加载局部地图，避免整图常驻内存带来的资源浪费。对大尺度矿井巷道场景而言，该架构能够在保证地图先验有效性的同时，增强系统的运行效率与工程可部署性。

\section{高精地图瓦片化分区与元数据组织}

\subsection{高精地图瓦片化分区方法}

为适应矿井巷道大尺度、长距离和多分支的场景特点，本文采用瓦片化方式对高精地图进行分区组织。具体而言，以二维平面坐标系为基准，按照固定的空间分辨率分别沿 $x$ 方向和 $y$ 方向对全局高精地图进行规则划分，形成若干局部地图瓦片。每个瓦片保存对应区域内的点云数据，瓦片尺寸由两个方向上的分区分辨率共同决定。该方式能够将原始整图拆分为多个空间上相互独立、逻辑上可组合调用的局部子图，从而为动态加载提供数据基础。高精地图瓦片划分示意如图~\ref{fig:ch5_tile_partition} 所示。

\begin{figure}[htbp]
    \centering
    \fbox{
        \parbox{0.88\linewidth}{
            \centering
            图5.2 高精地图瓦片化分区示意图占位图\\[4pt]
            可绘制矿井巷道平面地图，并用规则网格对地图进行切分，标注若干局部瓦片编号以及当前定位区域附近的邻域瓦片。
        }
    }
    \caption{高精地图瓦片化分区示意图}
    \label{fig:ch5_tile_partition}
\end{figure}

与直接存储单个整图文件相比，瓦片化组织具有明显优势。一方面，整图被拆分为多个局部瓦片后，系统可以按需读取目标区域对应的点云数据，降低单次地图调用规模；另一方面，瓦片划分后更有利于后续索引管理和局部区域维护，有助于高精地图在大场景下的组织与扩展。

\subsection{地图元数据索引组织}

为支撑在线快速检索，瓦片化分区之后还需建立配套的元数据索引。本文为每个瓦片记录其文件名及左下角起始坐标，同时在元数据文件中保存 $x$ 方向与 $y$ 方向的分区分辨率。对于任意瓦片 $M_i$，可将其索引信息表示为
\begin{equation}
M_i=\left\{f_i,o_i^x,o_i^y,r_x,r_y\right\}
\end{equation}
其中，$f_i$ 表示瓦片文件名，$o_i^x$ 和 $o_i^y$ 分别表示该瓦片左下角原点坐标，$r_x$ 和 $r_y$ 分别表示地图在两个方向上的分区分辨率。

通过元数据索引，系统在运行时无需遍历整幅地图点云内容，只需根据当前位置计算所属网格原点，并在索引表中快速找到相邻瓦片集合。相比直接按文件顺序逐一检索的方式，该方法能够显著提高地图查询效率，满足定位系统对在线调用速度的要求。

从论文整体逻辑上看，瓦片化分区与元数据索引并不改变第三章构建得到的高精地图本体，而是在其基础上完成结构化组织。因此，该部分内容与第三章离线高精地图构建方法并不冲突，而是对其在定位系统中的调用方式进行了补充和扩展。

\section{面向定位的局部地图动态加载与调度方法}

\subsection{动态加载基本思想}

在定位过程中，激光雷达在任一时刻仅能观测当前位置附近的局部环境，因此没有必要将全局高精地图完整加载到内存中。基于此，本文提出局部地图动态加载策略：根据当前位姿所在区域，从瓦片化高精地图中选取与当前位置相邻的一组局部瓦片，拼接形成当前时刻的局部匹配地图，并将其作为点云配准的先验地图。随着定位载体在巷道中运动，系统持续更新所需瓦片集合，实现局部地图的动态切换。该过程如图~\ref{fig:ch5_dynamic_loading} 所示。

\begin{figure}[htbp]
    \centering
    \fbox{
        \parbox{0.92\linewidth}{
            \centering
            图5.3 局部地图动态加载与调度流程图占位图\\[4pt]
            可按“当前位置估计 $\rightarrow$ 所属网格原点计算 $\rightarrow$ 邻域瓦片选择 $\rightarrow$ 瓦片请求发布 $\rightarrow$ 地图服务端加载与拼接 $\rightarrow$ 局部地图发布”的流程绘制。
        }
    }
    \caption{局部地图动态加载与调度流程图}
    \label{fig:ch5_dynamic_loading}
\end{figure}

与整图加载方式相比，动态加载策略使定位所需地图范围始终与当前位置保持一致，从而将匹配计算限制在空间邻域内，降低地图点云规模对内存和计算效率的影响。该方法本质上是高精地图的在线调度机制，而非地图本身的在线更新。

\subsection{基于当前位置的邻域瓦片选择方法}

设当前定位结果在地图坐标系下的位置为 $(x,y)$，瓦片分辨率分别为 $r_x$ 和 $r_y$，则当前位置所属网格原点可表示为
\begin{equation}
o_x=\left\lfloor \frac{x}{r_x}\right\rfloor r_x,\qquad
o_y=\left\lfloor \frac{y}{r_y}\right\rfloor r_y
\end{equation}
式中，$\lfloor \cdot \rfloor$ 表示向下取整运算。通过上述计算，可将当前位置映射到对应的地图网格单元中。

在得到当前位置所属网格原点后，进一步以该原点为中心，按照给定的邻域半径 $R$ 选取周围瓦片集合，构造当前时刻所需的局部地图。对应的瓦片集合可表示为
\begin{equation}
\mathcal{T}(o_x,o_y)=
\left\{
M_{ij}\mid i,j\in[-R,R]
\right\}
\end{equation}
其中，$M_{ij}$ 表示相对于当前原点在 $x$ 与 $y$ 方向偏移后的邻域瓦片。

邻域半径 $R$ 决定了局部地图覆盖范围。当 $R$ 取值较小时，系统加载地图范围较小，内存占用较低；当 $R$ 取值较大时，局部地图包含更多上下文信息，有利于增强边界区域的匹配稳定性。因此，邻域半径的选取需要在定位稳定性与实时性之间进行权衡。

\subsection{动态调度触发条件与切换稳定性控制}

若每一帧激光数据都重新计算并请求瓦片集合，将导致地图频繁切换，不仅增加磁盘读取和点云拼接开销，还可能在区域边界附近引入不必要的抖动。为此，本文在动态调度中引入距离门控和原点滞回两种稳定性控制策略。

首先，设置最小调度距离阈值 $d_{\mathrm{th}}$。仅当当前位置与上一次允许进行瓦片评估的位置之间的欧氏距离满足
\begin{equation}
\left\|
\mathbf{p}_k-\mathbf{p}_{k-1}^{\mathrm{eval}}
\right\|
\ge d_{\mathrm{th}}
\end{equation}
时，系统才重新计算所需瓦片集合。该策略避免了在短距离运动或局部抖动情况下频繁触发地图切换。

其次，为抑制载体在瓦片边界附近往返运动时出现的频繁切换现象，本文采用原点滞回机制。即当系统检测到当前位置对应的瓦片原点发生变化时，不立即切换到新原点，而是将其作为候选原点；只有当该候选原点在连续 $N_h$ 次有效评估中均保持一致时，系统才正式将稳定原点更新为新原点。该方法本质上引入了一个切换确认过程，有效提高了动态地图调度的稳定性。其示意如图~\ref{fig:ch5_switch_stability} 所示。

\begin{figure}[htbp]
    \centering
    \fbox{
        \parbox{0.86\linewidth}{
            \centering
            图5.4 地图切换稳定性控制示意图占位图\\[4pt]
            可绘制车辆在瓦片边界附近运动的示意图，并说明距离门控和原点滞回如何抑制频繁切换。
        }
    }
    \caption{地图切换稳定性控制示意图}
    \label{fig:ch5_switch_stability}
\end{figure}

通过距离门控与原点滞回策略的联合约束，系统能够在保持动态加载灵活性的同时，减少不必要的瓦片切换次数，从而提升定位过程中的地图调用稳定性。

\subsection{瓦片请求与局部地图生成流程}

在确定所需瓦片集合后，定位节点将对应的瓦片文件名列表封装为地图请求消息，并发送给地图服务模块。地图服务模块接收到请求后，依次读取目标瓦片文件，对其进行必要的下采样处理与点云拼接，生成当前时刻的局部地图，并以统一坐标系发布至定位模块。随后，定位模块利用当前激光帧与该局部地图进行配准，完成位姿估计。

该设计实现了定位节点与地图服务节点的功能分离，使地图调度、地图读取和定位计算能够相对独立地运行，有利于提高系统结构清晰度与工程实现稳定性。同时，由于局部地图由多个邻域瓦片动态拼接而成，其空间覆盖范围能够随载体运动而自适应变化，因此适用于大尺度矿井巷道环境下的连续定位任务。

\section{实验验证与结果分析}

\subsection{实验目的与设置}

为验证本文提出的高精地图瓦片化组织与动态加载方法的有效性，设计对比实验，将整图加载方案与动态加载方案进行比较，重点分析二者在内存占用、地图调度时延、配准耗时及定位连续性方面的差异。实验场景选取矿井巷道典型运行区域，地图数据采用第三章构建得到的高精地图，定位方法采用第四章所述定位流程。实验平台与参数设置如表~\ref{tab:ch5_exp_params} 所示。

在实验过程中，将整图加载方案作为基线方法，即在定位开始前一次性加载全部高精地图；动态加载方案则按照本章提出的方法，将高精地图划分为多个瓦片，并在定位过程中根据当前位置动态请求邻域瓦片。为进一步分析参数对系统性能的影响，分别调整邻域半径 $R$、最小调度距离阈值 $d_{\mathrm{th}}$ 以及滞回次数 $N_h$，比较不同配置下系统运行效果。实验场景及地图分区情况如图~\ref{fig:ch5_exp_scene} 所示。

\begin{table}[htbp]
    \centering
    \caption{实验平台与关键参数配置}
    \label{tab:ch5_exp_params}
    \small
    \renewcommand{\arraystretch}{1.15}
    \begin{tabular}{p{3.2cm}p{4.2cm}p{5.0cm}}
        \toprule
        模块或参数 & 配置项 & 说明 \\
        \midrule
        地图数据 & 第三章优化后的高精地图 & 作为整图加载与动态加载的统一地图输入。 \\
        分区参数 & $r_x$、$r_y$ & 分别表示地图在 $x$、$y$ 方向上的瓦片分辨率，具体数值待补充。 \\
        邻域范围参数 & $R$ & 控制当前位置周围参与拼接的邻域瓦片范围。 \\
        调度稳定参数 & $d_{\mathrm{th}}$、$N_h$ & 分别对应最小调度距离阈值与原点滞回次数。 \\
        定位模块 & 第四章定位方法 & 采用统一定位方法，保证实验比较公平。 \\
        实验平台 & 待补充 & 可补充处理器型号、内存大小及软件环境。 \\
        \bottomrule
    \end{tabular}
\end{table}

\begin{figure}[htbp]
    \centering
    \fbox{
        \parbox{0.90\linewidth}{
            \centering
            图5.5 实验场景与地图分区示意图占位图\\[4pt]
            可放置实验巷道地图、车辆轨迹及瓦片划分结果，并标注整图加载与局部动态加载的对比区域。
        }
    }
    \caption{实验场景与地图分区示意图}
    \label{fig:ch5_exp_scene}
\end{figure}

\subsection{评价指标}

为全面衡量所提方法的性能，本文采用以下指标进行评价。

（1）峰值内存占用：反映定位过程中地图加载对系统内存资源的消耗情况。

（2）平均地图调度时延：指从触发地图请求到局部地图生成完成的平均耗时，用于衡量动态地图调用效率。

（3）平均配准耗时：指单帧点云与局部地图进行匹配计算所需的平均时间，用于衡量定位实时性。

（4）地图切换次数：统计定位过程中局部地图实际发生更新的次数，用于分析动态调度策略的稳定性。

（5）定位连续性：通过轨迹连续程度、定位中断情况以及边界切换时的匹配稳定性，评价动态地图切换对定位结果的影响。

\subsection{实验结果与分析}

整图加载方案与动态加载方案的实验结果如表~\ref{tab:ch5_result_compare} 所示。从结果可以看出，动态加载方案在峰值内存占用方面明显优于整图加载方案，说明将高精地图划分为局部瓦片并按需调用，能够有效降低大场景地图常驻内存带来的资源压力。对于矿井巷道这类长距离场景，该优势尤为明显。

在平均配准耗时方面，动态加载方案同样表现出更好的实时性。这是因为整图加载时匹配模块需要面对更大规模的先验地图数据，而动态加载将配准对象限制在当前位置附近的局部地图范围内，从而降低了单次匹配计算复杂度。随着地图规模增大，动态加载在时间效率上的优势将更加突出。

在地图切换稳定性方面，若不引入距离门控与原点滞回机制，系统在瓦片边界附近易出现频繁切换，导致局部地图更新次数增加，并对匹配过程造成干扰。加入稳定性控制后，地图请求次数明显减少，局部地图切换更加平滑，定位轨迹连续性得到改善。由此可见，距离门控与原点滞回策略对于保证动态调度稳定性是必要的。

此外，不同参数配置下的对比结果如图~\ref{fig:ch5_param_analysis} 所示。当邻域半径较小时，系统内存占用和匹配耗时较低，但边界区域可用地图范围较小，容易影响匹配稳定性；当邻域半径适当增大后，局部地图上下文更加完整，定位连续性更好，但会带来一定的内存与计算开销。因此，在实际应用中应根据矿井巷道尺度与设备算力合理选择邻域半径。

\begin{table}[htbp]
    \centering
    \caption{整图加载与动态加载方案对比结果}
    \label{tab:ch5_result_compare}
    \small
    \renewcommand{\arraystretch}{1.15}
    \begin{tabular}{p{3.2cm}p{2.5cm}p{2.5cm}p{2.5cm}}
        \toprule
        指标 & 整图加载 & 动态加载 & 说明 \\
        \midrule
        峰值内存占用 / MB & 待补充 & 待补充 & 反映地图常驻内存压力。 \\
        平均地图调度时延 / ms & 不适用或待补充 & 待补充 & 动态加载方案需统计请求到局部地图生成的平均时延。 \\
        平均配准耗时 / ms & 待补充 & 待补充 & 比较两种地图组织方式下的单帧匹配效率。 \\
        地图切换次数 & 待补充 & 待补充 & 反映动态调度过程的稳定性。 \\
        定位连续性表现 & 待补充 & 待补充 & 可结合轨迹连续性或异常中断情况描述。 \\
        \bottomrule
    \end{tabular}
\end{table}

\begin{figure}[htbp]
    \centering
    \fbox{
        \parbox{0.90\linewidth}{
            \centering
            图5.6 参数影响分析图占位图\\[4pt]
            可绘制不同 $R$、$d_{\mathrm{th}}$ 和 $N_h$ 参数下的内存占用、匹配耗时或地图切换次数对比曲线。
        }
    }
    \caption{不同参数配置下的系统性能对比示意图}
    \label{fig:ch5_param_analysis}
\end{figure}

为便于后续补充真实实验数据，本章实验结果也可直接按如下描述进行替换：在实验场景中，整图加载方案的峰值内存占用为待补充数值，平均配准耗时为待补充数值；采用本文提出的动态加载方案后，峰值内存占用降低至待补充数值，平均配准耗时降低至待补充数值。与此同时，地图切换次数由待补充数值变化为待补充数值，边界区域定位轨迹更加连续。结果表明，所提高精地图瓦片化组织与动态加载方法能够在保证定位连续性的前提下，有效降低系统资源开销，提高大尺度矿井巷道场景下的定位实时性与工程可部署性。

\section{本章小结}

本章针对大尺度矿井巷道场景下高精地图整图加载开销大、定位实时性不足的问题，提出了面向矿井巷道定位的高精地图瓦片化组织与动态加载方法。首先，对离线优化后的高精地图进行规则瓦片化分区，并构建包含瓦片空间坐标与分辨率信息的元数据索引；其次，根据当前位姿估计结果设计邻域瓦片选择方法，实现局部地图的按需加载；进一步，引入距离门控和原点滞回机制，提升地图切换过程的稳定性；最后，通过整图加载与动态加载对比实验验证了所提方法在降低内存占用、缩短匹配耗时以及增强定位连续性方面的有效性。

本章研究表明，高精地图不仅需要具备较高的几何精度，还需要通过合理的组织与调度方式服务于定位系统运行。所提出的方法在不改变第三章高精地图构建原则和第四章定位方法核心流程的基础上，提升了系统在大场景矿井环境中的工程应用能力，为矿井巷道定位系统的实际部署提供了支撑。

\newpage
